\chapter{Related Work}
\label{ch:related-work}

\section{Federated Learning for IoT}
\label{sec:fl-iot}
% FL on resource-constrained devices; Ramadan et al. (2025)
% Challenges: communication cost, device heterogeneity, data heterogeneity

\section{Federated Learning over LoRaWAN: Torres Sanchez et al.}
\label{sec:torres-sanchez}

Torres Sanchez et al.~\cite{torressanchez2024federated} presented the first comprehensive study of \gls{fl} within \gls{lorawan}-constrained communication, targeting anomaly detection for industrial machinery.

\subsection{Their Experimental Setup}
% 4 IIoT prototypes: Manitou, Atlas Copco drill, Jaw Crusher, Doosan loader
% Private LoRaWAN → TTN → MQTT; 77,798 messages
% Flower framework, 4 virtual clients, 1 central server
% Model: autoencoder (32 neurons), Tanh, Adam, MSE loss
% Simulation environment: i7-1165G7, 16GB RAM, WSL + Conda

\subsection{Their Key Results}
% FL: Acc=92.30%, F1=94.77%, TNR=90.65%, TPR=92.93%
% Comparable to centralized OC-SVM (92.06%), AE (93.13%)
% Best config: 20 epochs × 4 rounds → F1=95.23%
% Model transfer: 1.39 KB/round at 32 neurons; 44 KB at 3 hidden layers

\subsection{Their Communication Analysis}
% Nm = ceil(Msize/MaxPayload) × Rd
% 1.39 KB model: 7 msgs (SF7) to 28 msgs (SF12) per round
% Training airtime: 52.8 minutes optimal

\section{ML-Based Path Loss Prediction}
\label{sec:ml-pathloss-related}
% ANN, RF, CNN for path loss; indoor vs outdoor
% Environment-driven approaches using weather/occupancy data
% Gap: NO federated approach for path loss prediction exists

\section{Research Gap and Positioning}
\label{sec:research-gap}

\begin{table}[H]
  \centering
  \begin{tabular}{p{3.5cm} p{4cm} p{4cm}}
    \toprule
    \textbf{Aspect} & \textbf{Torres Sanchez et al.} & \textbf{This Thesis} \\
    \midrule
    Application       & Anomaly detection (machinery) & Path loss regression (indoor LoRaWAN) \\[3pt]
    ML task           & Classification / reconstruction & Regression ($R^2$, RMSE) \\[3pt]
    Model             & Autoencoder (32 neurons) & Dense NN ($8\to8\to1$, 81 params) \\[3pt]
    FL framework      & Flower (Python) & Custom FedAvg (TF/Keras) \\[3pt]
    Clients           & 4 virtual & 6 virtual (one per real device) \\[3pt]
    Data              & 77,798 machinery msgs & 1,715,869 real LoRaWAN measurements \\[3pt]
    Non-IID split     & Per machine type & Per physical room location \\[3pt]
    HW prototype      & LoRaWAN prototype & TFLite Micro on real MKR WAN 1310 \\[3pt]
    Prior baseline    & OC-SVM, IF, AE & XGBoost $R^2 \approx 0.93$ \\[3pt]
    Update size       & 1.39\,KB (7+ msgs) & 52\,B (\textbf{1 msg}) \\
    \bottomrule
  \end{tabular}
  \caption{Positioning: Torres Sanchez et al.\ vs.\ this thesis.}
  \label{tab:positioning}
\end{table}

% "Torres Sanchez et al. proved FL is feasible over LoRaWAN.
%  Our prior project proved environment drives path loss.
%  This thesis bridges both: federated path loss regression on real edge hardware."


% ══════════════════════════════════════════════════════════════════════════════
% CHAPTER 4: METHODOLOGY
% ══════════════════════════════════════════════════════════════════════════════
